\subsection{A bounded answer to the Weyl-field automorphism problem}
\label{dev:weyl}

\paragraph{Question, scope and evidence.}
Item 1965-80.10 asks, literally, \ru{Найти все автоморфизмы этого тела}.
The strongest formulation supported by its preceding notation is to determine
\(\operatorname{Aut}_{L}D_{n,k}\), once a characteristic-zero coefficient field
\(L\) and its fixedness are specified. The page supplies neither proposed
generators nor relations for that group. In particular, its parenthesis about
Fourier transformation and multiplication is a suggestion, not a generation
conjecture. The following work fixes
\[
 K=\mathbb C(x),\qquad R=K[p;d/dx],\qquad
 D=Q(R)=D_{1,0},\qquad [p,x]=1.
\]
It determines the subgroup preserving this particular \(K\), and distinguishes
its rational translations that are inner from those that are outer. All
automorphisms fix \(\mathbb C\) pointwise. These hypotheses are restrictions
made here; they are not attributed to the notebook.

The full polynomial-algebra automorphism theorem is different:
\citet[Theorem 2]{makarlimanov1984automorphisms} treats \(A_1\), not its entire
fraction field. The rational-subfield stabilizer below is already
\citet[Proposition 2.2(i)]{alevdumas1998invariants}; their Theorem 2.5 also gives
\(D^G\simeq D\) for finite groups stabilizing \(\mathbb C(p)\), \(\mathbb C(x)\)
or \(\mathbb C(px)\). We re-prove the stabilizer and derive an explicit
inner-translation criterion. These are bounded, self-contained partial
answers, with no claim of new literature results. The notation \ours\ means
that a proof is supplied in this companion, not that priority is claimed.

\citet[\S5.3.3]{dumas2011invariants} explicitly records the missing full
automorphism-group description. The current search, dated 9 November 2026,
found no later full classification. A recent related result,
\citet[Theorem 0.12]{schwarz2024generalizations}, proves rationality of invariant
Weyl fields for finite groups induced from rational function fields when the
commutative invariant field is rational. It does not classify arbitrary
automorphisms of \(D\). A still more recent preprint,
\citet[Corollary 8.12]{lauritzenthomsen2026duality}, version 2 dated
2 November 2026, proves equality and finiteness of the left and right
indices for every noncommutative division \(\mathbb C\)-subalgebra of
\(D_{1,0}\); commutative subalgebras have infinite indices. This concerns
extension dimensions, not a description of all automorphisms.
These distinctions identify the remaining mathematical obstruction,
but do not establish why particular researchers did or did not pursue the
1966 suggestion; no documentary claim of that kind is made.

\paragraph{A centralizer calculation.}
We use formal pseudodifferential series only as an algebraic device. Let
\(\mathcal P\) consist of sums \(\sum_{j\le m}a_j(x)p^j\), with
\(a_j\in K\), \(m\in\mathbb Z\), and multiplication determined by
\[
 p^j a=\sum_{r\ge0}\binom{j}{r}a^{(r)}p^{j-r}
 \quad(j\in\mathbb Z).
\]
Every coefficient of a product receives finitely many contributions. For
completeness, associativity reduces, after the Leibniz rule, to the identity
\[
 \binom{i}{r}\binom{i+j-r}{s}
 =\sum_{t=0}^{s}\binom{j}{t}\binom{i}{r+s-t}
                         \binom{r+s-t}{r}
 \quad(i,j\in\mathbb Z,\ r,s\ge0).
\]
Indeed, first take out \(\binom{i}{r}\) and then apply the generalized
Vandermonde identity to the remaining sum. The two sides are exactly the
coefficients of \(a b^{(r)}c^{(s)}p^{i+j+k-r-s}\) in the two bracketings
of \((ap^i)(bp^j)(cp^k)\). Coefficientwise finiteness extends this calculation
from monomials to all the stated series.
A nonzero series can be factored as \(a_mp^m(1+Q)\), where \(Q\) has order
at most \(-1\); the formal geometric series in \(-Q\) constructs its inverse.
Thus \(\mathcal P\) is a division ring containing \(R\), and the Ore
universal property embeds \(D\) into it. No analytic convergence is asserted.

\begin{lemma}[The rational centralizer \ours]
The centralizer of \(x\) in \(D\) is precisely \(K\).
\end{lemma}
\begin{proof}
The multiplication rule gives \([p^j,x]=jp^{j-1}\) for every integer \(j\),
including negative integers, since all derivatives of \(x\) above the first
vanish. If \(A=\sum_{j\le m}a_jp^j\) commutes with \(x\), uniqueness of its
series expansion gives \(ja_j=0\) for every \(j\ne0\). Characteristic zero
therefore gives \(A=a_0\in K\). Conversely, every element of \(K\) commutes
with \(x\).
\end{proof}

\begin{proposition}[The rational-subfield stabilizer \ours]
An automorphism \(\phi\in\operatorname{Aut}_{\mathbb C}D\) satisfies
\(\phi(K)=K\) if and only if
\[
 \phi(x)=f(x)=\frac{ax+b}{cx+d},\qquad
 \phi(p)=\frac{1}{f'(x)}p+h(x),
 \quad ad-bc\ne0,\quad h\in K.
\]
Scalar multiples of the matrix represent the same \(f\). Every displayed pair
defines an automorphism, including when \(c\ne0\).
\end{proposition}
\begin{proof}
The restriction to \(K\) is a rational function field automorphism. Its image
\(f\) has a rational compositional inverse. Degrees of nonconstant rational
maps of the projective line multiply under composition, so \(\deg f=1\):
\(f\) is the displayed fractional linear map. In particular \(f'\ne0\).
The defining relation gives \([\phi(p),f]=1\), while
\([f'^{-1}p,f]=1\). Their difference commutes with \(f\), hence with
\(\mathbb C(f)=K\). The lemma puts that difference in \(K\).

Conversely, substitution \(a(x)\mapsto a(f(x))\) is an automorphism of \(K\).
For \(P=f'^{-1}p+h\), the chain rule gives
\([P,a(f)]=a'(f)\); thus the prescription respects the Ore relation and
defines a map on \(R\). Its images generate \(R\), since \(K\) is already
in the image and \(p=f'(P-h)\). More explicitly, if \(u=f^{-1}\), its inverse
has \(x\mapsto u(x)\) and
\[
 p\mapsto f'(u(x))p-f'(u(x))h(u(x)).
\]
The chain rule verifies both compositions on the generators. The maps extend
to inverse automorphisms of \(D\).
\end{proof}

The equality \(\phi(K)=K\) is indispensable: the algebraic Fourier
automorphism \(x\mapsto p\), \(p\mapsto-x\) does not preserve \(K\).
Conversely, preserving \(K\) need not mean preserving \(A_1\):
\(x\mapsto x^{-1}\), \(p\mapsto-x^2p\) is an involution of \(D\) whose image
of \(x\) is not polynomial. This is a statement about the chosen polynomial
subalgebra, not about whether the involution is conjugate to another action.

\paragraph{Which rational translations are inner?}
Write \(T_h(x)=x\) and \(T_h(p)=p+h(x)\).
The maps \(T_h\) form a copy of the additive group of \(K\).

\begin{proposition}[Exact rational-pole criterion \ours]
For \(h\in\mathbb C(x)\), the following are equivalent:
\begin{enumerate}[label=(\roman*)]
\item \(T_h\) is an inner automorphism of \(D\).
\item \(h=-u'/u\) for some \(u\in\mathbb C(x)^{\times}\).
\item The polynomial part of \(h\) is zero, every finite pole is simple,
and every finite-pole residue is an integer.
\end{enumerate}
If \(h=\sum_a r_a/(x-a)\) has these properties, one implementing element is
\(u=\prod_a(x-a)^{-r_a}\).
\end{proposition}
\begin{proof}
If \(T_h=\operatorname{Ad}_u\) for some \(u\in D^{\times}\), its fixing
\(x\) implies \(ux=xu\). The centralizer lemma forces \(u\in K^{\times}\).
The Ore relation yields
\[
 upu^{-1}=p-\frac{u'}u.
\]
This proves (i) implies (ii), and the same identity proves the converse.
Factor a nonzero rational function over \(\mathbb C\) as
\(u=C\prod_a(x-a)^{m_a}\), with integral \(m_a\) and finitely many nonzero
exponents. Then \(-u'/u=-\sum_a m_a/(x-a)\), proving (ii) implies (iii).
Conversely, unique partial fraction decomposition and (iii) give exactly
\(h=\sum_a r_a/(x-a)\), \(r_a\in\mathbb Z\); the displayed product has
\(-u'/u=h\). No extra condition at infinity is missing: the differential
\(h(x)\,dx\) has residue there equal to the negative sum of the finite
residues.
\end{proof}

\begin{corollary}[Translation classes in the outer automorphism group \ours]
The image of \(\{T_h:h\in K\}\) in
\(\operatorname{Out}_{\mathbb C}D=\operatorname{Aut}_{\mathbb C}D/\operatorname{Inn}D\)
is, as an abelian group,
\[
 K\Big/\left\{-\frac{u'}u:u\in K^{\times}\right\}
 \simeq
 \mathbb C[x]\ \oplus\!\bigoplus_{a\in\mathbb C,\ j\ge2}\!\mathbb C
 \ \oplus\!\bigoplus_{a\in\mathbb C}\!(\mathbb C/\mathbb Z).
\]
All direct-sum elements have finite support. The first term records the
polynomial part, the second higher-order principal parts, and the third
simple-pole residues modulo integers.
\end{corollary}
\begin{proof}
The kernel of the homomorphism \(h\mapsto[T_h]\) is exactly the subgroup in
the denominator by the proposition. Partial fractions give an additive direct
sum decomposition of \(K\). The denominator is precisely the integral lattice
in its simple-pole summands, so quotienting gives the formula. This is an
abelian-group quotient; the \(\mathbb C/\mathbb Z\) terms are not complex
vector spaces.
\end{proof}

\begin{example}[Outer torsion and actual infinite order \ours]
For an integer \(m\ge2\), \(T_{1/(mx)}\) is not inner, but
\(T_{1/(mx)}^m=T_{1/x}=\operatorname{Ad}_{x^{-1}}\). Its class in the outer
group has exact order \(m\), because \(r/m\notin\mathbb Z\) for
\(1\le r<m\). The automorphism itself has infinite order:
\(T_{1/(mx)}^r=T_{r/(mx)}\ne\mathrm{id}\) for every positive \(r\).
Thus statements about finite order in \(\operatorname{Out}D\) must not be
substituted for statements about finite order in \(\operatorname{Aut}D\).
\end{example}

\paragraph{Finite subgroups of this stabilizer.}
For \(\sigma\in\operatorname{Aut}_{\mathbb C}K\) put \(f=\sigma(x)\) and
let \(L_\sigma\) be the lift with \(p\mapsto f'^{-1}p\).
Then
\[
 L_\sigma T_b L_\sigma^{-1}=T_{R_\sigma b},\qquad
 R_\sigma b=f'(x)b(f(x)).
\]
With automorphisms composed as maps of algebras, \(f_{\sigma\tau}=
f_\tau\circ f_\sigma\), and the chain rule gives
\(R_{\sigma\tau}=R_\sigma R_\tau\). Thus every stabilizer element has a
unique form \(T_bL_\sigma\), with multiplication
\[
 (T_bL_\sigma)(T_cL_\tau)=T_{b+R_\sigma c}L_{\sigma\tau}.
\]

\begin{proposition}[Removing a finite-group translation cocycle \ours]
Every finite subgroup \(H\) of the stabilizer of \(K\) is conjugate, by a
rational translation \(T_a\), to the canonical lift of its image in
\(\operatorname{Aut}_{\mathbb C}K\).
\end{proposition}
\begin{proof}
The kernel of restriction to \(K\) is the torsion-free additive group
\(\{T_b\}\), so restriction is injective on \(H\). Write its image as
\(Q\), and its elements uniquely as \(T_{b_\sigma}L_\sigma\),
\(\sigma\in Q\). The multiplication law gives
\(b_{\sigma\tau}=b_\sigma+R_\sigma b_\tau\).
Let \(a=|Q|^{-1}\sum_{\tau\in Q}b_\tau\). Sum this identity over \(\tau\)
to obtain \(b_\sigma=a-R_\sigma a\). Hence
\[
 T_{-a}(T_{b_\sigma}L_\sigma)T_a=L_\sigma
\]
for every \(\sigma\). Averaging uses characteristic zero; the assertion is
not made in positive characteristic.
\end{proof}

This proves a normal form for finite actions within the chosen stabilizer.
It says nothing about whether an arbitrary finite subgroup of
\(\operatorname{Aut}D\) stabilizes a conjugate rational subfield. Nor does
it make the conjugating translation inner: the pole criterion decides that
additional question separately.

\paragraph{Computations and independent checks.}
The accompanying \texttt{code/check\_weyl.py} implements multiplication of
finite differential operators from the Leibniz rule, with exact rational
coefficients. It checks the commutator, inverse and composition identities
above for fractional linear changes including inversion, higher-order poles,
nonintegral residues and polynomial shifts. A separate evaluation on
monomials checks the operator commutator without using the multiplication
routine. The same script enumerates PBW exponent tuples by a recurrence and
compares them with \(\binom{N+2n+k}{2n+k}\), including zero generators and
\(N=0\). The complete output is retained; finite tests are used to detect
algebraic or coding mistakes, not to establish the universally quantified
results.

\begin{example}[Rational shifts checked by symbolic differentiation \oursComputed]
The exact examples give
\[
\begin{array}{c|c|c}
h(x)&\text{inner?}&\text{implementer, when it exists}\\\hline
3/(x-1)-2/(x+2)&\text{yes}&(x-1)^{-3}(x+2)^2\\
1/(2x)&\text{no}&\text{its square is implemented by }x^{-1}\\
1/x^2&\text{no}&\text{higher-order pole}\\
x&\text{no}&\text{nonzero polynomial part}\\
0&\text{yes}&1
\end{array}
\]
Each positive case is checked by differentiating the implementer; negative
cases are excluded by the proved criterion, not by a failed search.
\end{example}

The second proof check re-derives the conjugation sign from
\(pu=up+u'\), and the composition law from substitution before applying the
abstract semidirect-product notation. Deliberately reversing the sign of
\(u'/u\), or omitting \(f'\) in the group action, fails the symbolic checks.
The proof also distinguishes the fixed subfield condition from the weaker
condition that an image is rational in some other generator. A separate
review of the written proof and a second parameter run are documented in
\texttt{evidence/development\_review.md} and the computation output.

\paragraph{The exact obstruction and a further bounded run.}
The unresolved step is global: an arbitrary automorphism can move
\(\mathbb C(x)\) to a different commutative subfield of \(D\). The
centralizer argument classifies maps only after the source and target
rational subfield have been fixed. Fourier interchange already escapes that
stabilizer. No reduction of every Weyl generating pair to the normal form
above has been proved here, and no generation theorem is conjectured on the
strength of the computations.

A further bounded run could explicitly compare the stabilizers of
\(\mathbb C(x)\) and \(\mathbb C(xp)\), using the second normal form in
\citet[Proposition 2.2(ii)]{alevdumas1998invariants}, and compute innerness
and conjugacy obstructions for rational functions with a fixed finite pole
set. That would refine the notebook's concrete suggestions while remaining
far short of classifying the full automorphism group. The present contribution
is the explicit, checked restricted answer and the precise location of the
remaining gap.
